\unit{Exposición 195 -- Técnica de descenso II}
\sid{195}
\src{369}
\seminar{12.\textsuperscript{o} año, 1959/60, n.\textsuperscript{o} 195\\Febrero de 1960}
\paperhead{TÉCNICA DE DESCENSO Y TEOREMAS DE EXISTENCIA EN GEOMETRÍA ALGEBRAICA\\
II. EL TEOREMA DE EXISTENCIA EN LA TEORÍA FORMAL DE LOS MÓDULOS}%
  {por Alexander GROTHENDIECK}

\fsection{A}{Funtores representables y pro-representables}

\subsection*{1. Funtores representables}
\addcontentsline{toc}{subsection}{1. Funtores representables}

Sea $\mathcal C$ una categoría. Para todo $X\in\mathcal C$, sea $h_X$ el
funtor contravariante de $\mathcal C$ en la categoría $(\operatorname{Ens})$
de los conjuntos,
\[
  h_X:\mathcal C\longrightarrow(\operatorname{Ens}),
\]
definido por la fórmula
\[
  h_X(Y)=\operatorname{Hom}(Y,X).
\]
Un morfismo $X\longrightarrow X'$ en $\mathcal C$ induce de manera evidente un
homomorfismo de funtores $h_X\longrightarrow h_{X'}$: $h_X$ es un funtor
covariante en $X$; es decir, hemos definido un funtor covariante canónico
\[
  h:\mathcal C\longrightarrow
  \operatorname{Hom}\bigl(\mathcal C^o,(\operatorname{Ens})\bigr)
\]
de $\mathcal C$ en la categoría de los funtores covariantes de la dual
$\mathcal C^o$ de $\mathcal C$ en la categoría de los conjuntos. Recordemos
entonces:

\result{PROPOSICIÓN}{1.1} Este funtor $h$ es \emph{plenamente fiel}; en otros
términos, para todo par $X,X'$ de objetos de $\mathcal C$, la aplicación
natural
\[
  \operatorname{Hom}(X,X')\longrightarrow
  \operatorname{Hom}(h_X,h_{X'})
\]
es biyectiva.

En particular, si un funtor
$F\in\operatorname{Hom}(\mathcal C^o,(\operatorname{Ens}))$ es isomorfo a un
funtor de la forma $h_X$, entonces $X$ queda determinado salvo isomorfismo
único. Se dice en tal caso que el funtor $F$ es \emph{representable}. La
proposición anterior significa que el funtor canónico $h$ define una
equivalencia de la categoría $\mathcal C$ con la subcategoría plena de
$\operatorname{Hom}(\mathcal C^o,(\operatorname{Ens}))$ formada por los
funtores representables. Este hecho fundamenta la noción de «solución de un
problema universal»: tal problema consiste siempre en examinar si un funtor
dado (contravariante como aquí, o covariante en el caso dual) de $\mathcal C$
en $(\operatorname{Ens})$ es representable.

\src{370}
Observemos además que, por la propia definición de producto en una categoría
[1], el funtor $h:X\longmapsto h_X$ conmuta con los productos siempre que estos
existan (y, más generalmente, con los límites proyectivos finitos o infinitos;
en particular, con los productos fibrados, con la formación de «núcleos» [2],
etc., siempre que existan): se tiene un isomorfismo de funtores
\[
  h_{X\times X'}\xrightarrow{\sim}h_X\times h_{X'}
\]
siempre que exista $X\times X'$; es decir, se tienen biyecciones functoriales
en $Y$
\[
  h_{X\times X'}(Y)\xrightarrow{\sim}h_X(Y)\times h_{X'}(Y).
\]
En particular, dar un morfismo
\[
  X\times X'\longrightarrow X''
\]
en $\mathcal C$ (es decir, una «ley de composición» en $\mathcal C$ entre
$X,X',X''$) equivale a dar un morfismo
\[
  h_{X\times X'}=h_X\times h_{X'}\longrightarrow h_{X''},
\]
es decir, a dar, para todo $Y\in\mathcal C$, una ley de composición de
conjuntos
\[
  h_X(Y)\times h_{X'}(Y)\longrightarrow h_{X''}(Y)
\]
de tal modo que, para todo morfismo $Y\longrightarrow Y'$ en $\mathcal C$, el
sistema de aplicaciones entre conjuntos
\[
  h_{X^{(i)}}(Y)\longleftarrow h_{X^{(i)}}(Y')
  \qquad(\text{para }i=0,1,2)
\]
sea un morfismo respecto de las dos leyes de composición correspondientes a
$Y$ y $Y'$. Se ve así que la noción de estructura de «$\mathcal C$-grupo»,
«$\mathcal C$-anillo», etc. sobre un objeto $X$ de $\mathcal C$ se expresa del
modo más cómodo (tanto en la teoría como en la práctica) diciendo que, para
todo $Y\in\mathcal C$, el conjunto $h_X(Y)$ está provisto de una ley de grupo
(resp. de anillo, etc.) en el sentido usual, y que las aplicaciones
$h_X(Y)\longleftarrow h_X(Y')$ correspondientes a morfismos
$Y\longrightarrow Y'$ son homomorfismos de grupos (resp. de anillos, etc.).
Esta es, por ejemplo, la forma más intuitiva y cómoda de definir los diversos
grupos clásicos $G_a$, $G_m$, $\operatorname{Gl}(n)$, etc. sobre un preesquema
base arbitrario $S$ y de escribir las relaciones clásicas entre esos grupos;
de dotar al esquema afín $V(\mathcal F)$ sobre $S$, definido por un haz
cuasicoherente $\mathcal F$, de una estructura de «fibrado vectorial»; y de
definir y estudiar las diversas variedades de banderas asociadas
(grassmannianas, fibrados proyectivos), etc. El yoga general consiste
sencillamente en identificar, mediante el funtor canónico $h$, los objetos de
$\mathcal C$ con ciertos funtores contravariantes: los funtores representables,

\src{371}
de $\mathcal C$ en la categoría de los conjuntos.

El procedimiento habitual de invertir las flechas, necesario por ejemplo en el
caso de los esquemas afines para pasar del lenguaje geométrico al del álgebra
conmutativa, conduce a dualizar las consideraciones anteriores y, en
particular, a introducir también los funtores covariantes representables
$\mathcal C\longrightarrow(\operatorname{Ens})$; es decir, los de la forma
\[
  Y\longmapsto\operatorname{Hom}(X,Y)=h'_X(Y).
\]

\subsection*{2. Funtores pro-representables, proobjetos}
\addcontentsline{toc}{subsection}{2. Funtores pro-representables, proobjetos}

Sea $\widetilde X=(X_i)_{i\in I}$ un sistema proyectivo de objetos de
$\mathcal C$; le corresponde un funtor covariante
\[
  h'_{\widetilde X}=\varinjlim_i h'_{X_i},
\]
más explícitamente,
\[
  h'_{\widetilde X}(Y)=\varinjlim_i h'_{X_i}(Y)
  =\varinjlim_i\operatorname{Hom}(X_i,Y),
\]
de $\mathcal C$ en $(\operatorname{Ens})$. Un funtor de $\mathcal C$ en
$(\operatorname{Ens})$ que sea isomorfo a un funtor de este tipo con $I$
filtrante se llama \emph{pro-representable}. Por el número anterior, estos son
exactamente los funtores isomorfos a límites inductivos filtrantes de funtores
representables. Sea $\widetilde X'=(X'_j)_{j\in J}$ un segundo sistema
proyectivo filtrante en $\mathcal C$ (construido sobre un segundo conjunto de
índices preordenado filtrante $J$); se comprueba fácilmente que hay una
biyección canónica (que generaliza la proposición 1.1)
\[
  \operatorname{Hom}(h'_{\widetilde X'},h'_{\widetilde X})
  =\varprojlim_j\varinjlim_i\operatorname{Hom}(X_i,X'_j).
\]
Esto lleva a introducir la categoría $\operatorname{Pro}(\mathcal C)$ de los
proobjetos de $\mathcal C$; sus objetos son los sistemas proyectivos de objetos
de $\mathcal C$ (sobre conjuntos variables de índices preordenados filtrantes),
y si $\widetilde X=(X_i)_{i\in I}$ y
$\widetilde X'=(X'_j)_{j\in J}$ son dos objetos de esta clase, se pone
\[
  \operatorname{ProHom}(\widetilde X,\widetilde X')
  =\varprojlim_j\varinjlim_i\operatorname{Hom}(X_i,X'_j),
\]
siendo evidente la composición de los prohomomorfismos. Por su propia
construcción, $\widetilde X\longmapsto h'_{\widetilde X}$ puede considerarse
como un funtor contravariante en $\widetilde X$, que establece una equivalencia
entre la categoría dual de la categoría $\operatorname{Pro}(\mathcal C)$ de
los proobjetos de $\mathcal C$ y la categoría de los funtores covariantes
pro-representables de $\mathcal C$ en $(\operatorname{Ens})$. Naturalmente, un
objeto $X$ de $\mathcal C$ define canónicamente un proobjeto, denotado

\src{372}
también por $X$, de modo que $\mathcal C$ es equivalente a una subcategoría
plena de $\operatorname{Pro}(\mathcal C)$. Si entonces
$\widetilde X=(X_i)_{i\in I}$ es un proobjeto cualquiera de $\mathcal C$, se
tendrá (con la identificación anterior)
\[
  \widetilde X=\varprojlim_i X_i,
\]
tomándose el límite proyectivo en $\operatorname{Pro}(\mathcal C)$
(puesto que $h'_{\widetilde X}=\varinjlim_i h'_{X_i}$). Hay que observar que,
aun cuando el límite proyectivo de los $X_i$ exista en $\mathcal C$, en
general no será isomorfo al límite proyectivo $\widetilde X$ en
$\operatorname{Pro}(\mathcal C)$, como ya resulta evidente cuando $\mathcal C$
es la categoría de los conjuntos. Nótese que, por definición,
$\varprojlim_{\mathcal C}X_i=L$ queda definido por la condición de que el
funtor en $Y\in\mathcal C$
\[
  \varprojlim_i\operatorname{Hom}_{\mathcal C}(Y,X_i)
  =\operatorname{Hom}_{\operatorname{Pro}(\mathcal C)}(Y,\widetilde X)
\]
con valores en $(\operatorname{Ens})$ sea representable mediante $L$, es decir,
sea isomorfo a $\operatorname{Hom}_{\mathcal C}(Y,L)$; por consiguiente,
$\varprojlim_{\mathcal C}X_i$ ya está definido en términos del proobjeto
$\widetilde X$ y, más precisamente, depende functorialmente del proobjeto
$\widetilde X$ siempre que esté definido; no hay, pues, inconveniente en
denotarlo por $\varprojlim_{\mathcal C}(\widetilde X)$. Si los límites
proyectivos existen siempre en $\mathcal C$,
$\varprojlim_{\mathcal C}(\widetilde X)$ es un funtor de
$\operatorname{Pro}(\mathcal C)$ en $\mathcal C$, y hay un homomorfismo
canónico de funtores
$\varprojlim_{\mathcal C}(\widetilde X)\longrightarrow\widetilde X$.

Como todo funtor, que supondremos covariante para fijar las ideas,
\[
  F:\mathcal C\longrightarrow\mathcal C'
\]
se prolonga de manera evidente a un funtor
\[
  \operatorname{Pro}(F):\operatorname{Pro}(\mathcal C)
  \longrightarrow\operatorname{Pro}(\mathcal C'),
\]
se sigue que, si en $\mathcal C'$ existen siempre los límites proyectivos, $F$
define también canónicamente un funtor compuesto
\[
  \overline F:\operatorname{Pro}(\mathcal C)\longrightarrow\mathcal C',
  \qquad
  \overline F=\varprojlim_{\mathcal C'}\circ\operatorname{Pro}(F),
\]
que transforma $\widetilde X=(X_i)_{i\in I}$ en
$\varprojlim_{\mathcal C'}F(X_i)$.

Un proobjeto $\widetilde X$ se llama \emph{proobjeto estricto} si es isomorfo a
un proobjeto $(X_i)_{i\in I}$ cuyos morfismos de transición
$X_i\longrightarrow X_j$ son epimorfismos; un funtor definido por tal objeto se
llama \emph{estrictamente pro-representable}. Se puede exigir además que $I$
sea ordenado filtrante y que todo epimorfismo $X_i\longrightarrow X'$ sea
equivalente a un epimorfismo

\src{373}
$X_i\longrightarrow X_j$ para un $j\in I$ adecuado (determinado de manera
única por esta condición). En estas condiciones, el sistema proyectivo
$(X_i)_{i\in I}$ queda determinado salvo isomorfismo único (en el sentido usual
de los isomorfismos de sistemas proyectivos). Se sigue que en
$\operatorname{Pro}(\mathcal C)$ existe siempre el límite proyectivo de un
sistema proyectivo de proobjetos estrictos $\widetilde X^{(\alpha)}$ y que, con
las notaciones anteriores para $F$, $\overline F$, se tendrá
\[
  \overline F\left(\varprojlim_\alpha\widetilde X^{(\alpha)}\right)
  =\varprojlim_{\alpha,\mathcal C'}
  \overline F\left(\widetilde X^{(\alpha)}\right).
\]
En particular, si todo proobjeto de $\mathcal C$ es estricto (cf. el número
siguiente), el funtor prolongado $\overline F$ conmuta con los límites
proyectivos.

\subsection*{3. Caracterización de los funtores pro-representables}
\addcontentsline{toc}{subsection}{3. Caracterización de los funtores pro-representables}

Sean $\mathcal C$, $\mathcal C'$ dos categorías en las que existen los límites
proyectivos finitos (es decir, relativos a conjuntos preordenados finitos, no
necesariamente filtrantes), o, lo que es equivalente, en las que existen los
productos finitos y los productos fibrados finitos (lo cual implica, en
particular, la existencia de un «objeto final» $e$, tal que
$\operatorname{Hom}(X,e)$ se reduzca a un elemento para todo $X$). Sea $F$ un
funtor covariante de $\mathcal C$ en $\mathcal C'$. Las condiciones siguientes
son equivalentes:

\begin{enumerate}
\item[(i)] $F$ conmuta con los límites proyectivos finitos;
\item[(ii)] $F$ conmuta con los productos finitos y los productos fibrados finitos;
\item[(iii)] $F$ conmuta con los productos finitos y, para todo diagrama exacto
\[
  X\longrightarrow X'\rightrightarrows X''
\]
en $\mathcal C$ ([3], A, definición 2.1), el diagrama transformado por $F$
\[
  F(X)\longrightarrow F(X')\rightrightarrows F(X'')
\]
es exacto.
\end{enumerate}

Se dice entonces que $F$ es \emph{exacto por la izquierda}.

En lo sucesivo se supone que en $\mathcal C$ existen los límites proyectivos
finitos. De las definiciones resulta entonces inmediato que un funtor
representable es exacto por la izquierda y, pasando al límite, que un funtor
pro-representable es exacto por la izquierda.

Para obtener un recíproco, sea
\[
  F:\mathcal C\longrightarrow(\operatorname{Ens})
\]
un funtor covariante, y sean $X\in\mathcal C$ y $\xi\in F(X)$. Se dice que
$\xi$ (o el par \src{374}$(X,\xi)$) es \emph{minimal} si, para todo
$X'\in\mathcal C$ y $\xi'\in F(X')$, y todo monomorfismo estricto ([3], A, 2)
$u:X'\longrightarrow X$ tal que $\xi=F(u)(\xi')$, $u$ es un isomorfismo. Por
otra parte, se dice que un par $(X,\xi)$ \emph{domina} a $(X'',\xi'')$
($\xi\in F(X)$, $\xi''\in F(X'')$) si existe un morfismo
$v:X\longrightarrow X''$ tal que $\xi''=F(v)(\xi)$; si $\xi$ es minimal y $F$
es exacto por la izquierda, este morfismo $v$ es entonces único; si $\xi''$ es
minimal, $v$ es sobreyectivo. Se deduce fácilmente la proposición siguiente:

\result{PROPOSICIÓN}{3.1} Para que $F$ sea estrictamente pro-representable, es
necesario y suficiente que satisfaga las dos condiciones:
\begin{enumerate}
\item[(i)] $F$ es exacto por la izquierda;
\item[(ii)] todo par $(X,\xi)$, con $\xi\in F(X)$, está dominado por un par
minimal.
\end{enumerate}

Esta última condición es vacía si todo objeto de $\mathcal C$ es artiniano
(tomando un subobjeto $X'$ de $X$ minimal entre aquellos para los cuales existe
$\xi'\in F(X')$ tal que $\xi$ sea la imagen de $\xi'$). De aquí:

\result{COROLARIO}{} Sea $\mathcal C$ una categoría cuyos objetos son
artinianos y en la que existen los límites proyectivos finitos. Entonces los
funtores pro-representables de $\mathcal C$ en $(\operatorname{Ens})$ son
exactamente los funtores exactos por la izquierda, e incluso son estrictamente
pro-representables.

Este último hecho significa también que \emph{todo proobjeto de $\mathcal C$
es estricto}.

\subsection*{4. Ejemplo: grupos de tipo galoisiano, grupos proalgebraicos}
\addcontentsline{toc}{subsection}{4. Grupos de tipo galoisiano, grupos proalgebraicos}

Si $\mathcal C$ es la categoría de los grupos finitos ordinarios,
$\operatorname{Pro}(\mathcal C)$ es equivalente a la categoría de los grupos
topológicos compactos totalmente disconexos. Son grupos de este tipo y sus
generalizaciones, obtenidas sustituyendo los grupos finitos ordinarios por
esquemas en grupos finitos sobre un preesquema base dado (por ejemplo, los
grupos algebraicos finitos sobre un cuerpo $k$), los que harán las veces de
grupos fundamentales, de homotopía y de homología absolutos y relativos para
los preesquemas. En todos estos ejemplos se aplica el corolario de la
proposición 3.1, y es precisamente mediante el funtor asociado como debe
definirse $\pi_1$ [2]. Lo mismo sucede partiendo de la categoría de los grupos
algebraicos o cuasialgebraicos sobre un cuerpo (o, más generalmente, sobre un
preesquema noetheriano): se obtienen los «grupos proalgebraicos» de Serre [4].

\subsection*{5. Ejemplo: «variedades formales»}
\addcontentsline{toc}{subsection}{5. Ejemplo: variedades formales}

Sean $\Lambda$ un anillo noetheriano y $\mathcal C$ la categoría de las
$\Lambda$-álgebras $A$ que son módulos artinianos de tipo finito sobre
$\Lambda$ (o, más brevemente,

\src{375}
las $\Lambda$-álgebras artinianas). Se cumplen las condiciones del corolario de
la proposición 3.1. La categoría $\operatorname{Pro}(\mathcal C)$ es aquí
equivalente a la categoría de las álgebras topológicas $O$ sobre $\Lambda$
isomorfas a límites proyectivos topológicos
\[
  O=\varprojlim_i O_i
\]
de álgebras $O_i\in\mathcal C$, es decir, cuya topología es lineal, separada y
completa, y tal que para todo ideal abierto $J_i$ de $O$, $O/J_i$ sea un
álgebra artiniana sobre $\Lambda$. El funtor
$\mathcal C\longrightarrow(\operatorname{Ens})$ asociado a tal álgebra no es
otro que
\begin{align*}
  F(A)=h'_O(A)
  &=\{\text{homomorfismos continuos de $\Lambda$-álgebras topológicas }
       O\longrightarrow A\}\\
  &=\varinjlim_i
    \operatorname{Hom}_{\Lambda\text{-álgebras}}(O_i,A).
\end{align*}

Obsérvese, además, que la categoría $\mathcal C$ es esencialmente el producto
de las categorías análogas correspondientes a las completaciones de los
anillos locales $\Lambda_{\mathfrak m}$, para los ideales maximales
$\mathfrak m$ de $\Lambda$; por tanto, si se desea, se puede restringir al caso
en que $\Lambda$ sea un anillo local completo de esta clase. En cualquier caso,
$O$ se descompone canónicamente en el producto topológico de componentes
locales, correspondientes a los «puntos» del esquema formal [2] definido por
$O$. Un punto de esta clase está definido por un elemento $\xi$ de algún
$F(K)$, donde $K\in\mathcal C$ es un cuerpo (por ejemplo, el cuerpo residual
de la componente local considerada); dos pares $(\xi,K)$ y $(\xi',K')$
definen el mismo punto si y solo si están dominados por un mismo
$(\xi'',K'')$, o bien si dominan a un mismo $(\xi''',K''')$. (Si los
$\Lambda/\mathfrak m$ son algebraicamente cerrados, basta tomar la suma
disjunta de los $F(\Lambda/\mathfrak m)$.)

Es importante dar condiciones para que la componente local $O_\xi$ de $O$
correspondiente a un $\xi\in F(K)$ sea un anillo noetheriano. Cuando $\Lambda$
es un anillo local completo (noetheriano, recordémoslo), esto equivale a decir
que $O_\xi$ es isomorfo a un anillo cociente de un anillo de series formales
$\Lambda[[t_1,\ldots,t_n]]$ sobre $\Lambda$. Para dar un criterio de esta clase,
introduzcamos (para todo anillo $A$) la $A$-álgebra $I_A$ de los «números
duales» de $A$ mediante
\[
  I_A=A[t]/t^2A[t].
\]
Sea $\varepsilon:I_A\longrightarrow A$ el homomorfismo de aumento; este define
(si $A\in\mathcal C$) una aplicación
\[
  F(\varepsilon):F(I_A)\longrightarrow F(A),
\]

\src{376}
y, utilizando que $F$ es exacto por la izquierda, se define de manera
intrínseca una estructura de $A$-módulo en la parte
\[
  F(I_A,\xi)=F(\varepsilon)^{-1}(\xi)
\]
de $F(I_A)$ formada por los $\xi'\in F(I_A)$ que «se reducen a $\xi$»;
utilizando la forma explícita de $F$ en términos de $O$, se encuentra que este
$A$-módulo se identifica con
\[
  \operatorname{Hom}_{\Lambda}(\mathfrak m_\xi/\mathfrak m_\xi^2,A),
\]
donde $\mathfrak m_\xi$ es el núcleo del homomorfismo
$\xi:O\longrightarrow A$, es decir, si $A$ es un cuerpo, el ideal maximal de
la componente local $O_\xi$ de $O$. Se deduce inmediatamente la proposición
siguiente:

\result{PROPOSICIÓN}{5.1} Sea $\xi\in F(K)$, donde $K\in\mathcal C$ es un
cuerpo. Para que la componente local correspondiente $O_\xi$ de $O$ sea un
anillo noetheriano, es necesario y suficiente que el conjunto $F(I_K,\xi)$ de
los elementos de $F(I_K)$ que se reducen a $\xi$ sea un espacio vectorial
de dimensión finita sobre $K$. En estas condiciones, se tiene un isomorfismo
canónico
\[
  F(I_K,\xi)=
  \operatorname{Hom}\!\left(
    \frac{\mathfrak m_\xi/\mathfrak m_\xi^2}
         {\operatorname{im}(\mathfrak n_\xi/\mathfrak n_\xi^2)},K
  \right),
\]
donde $\mathfrak n_\xi$ es el ideal maximal de $\Lambda$, núcleo del
homomorfismo $\Lambda\longrightarrow K$.
\footnote{El texto impreso da aquí una fórmula que solo es correcta cuando
$\Lambda$ es un cuerpo. La fe de erratas de 1962 prescribe, en el caso general,
el cociente de $\mathfrak m_\xi/\mathfrak m_\xi^2$ por la imagen de
$\mathfrak n_\xi/\mathfrak n_\xi^2$.}
En particular, la dimensión del $K$-espacio vectorial $F(I_K,\xi)$ es igual a
la dimensión del espacio vectorial así corregido sobre el cuerpo
$O_\xi/\mathfrak m_\xi=K(\xi)$.

Supongamos que $O_\xi$ es noetheriano y, para simplificar la escritura, que
$\Lambda$ es local completo y $O=O_\xi$. Se dice que $O$ es \emph{liso sobre
$\Lambda$} si $O$ es un álgebra finita y étale sobre el álgebra completada del
localizado de $\Lambda[t_1,\ldots,t_n]$ en un ideal maximal de este anillo que
induce el ideal maximal de $\Lambda$; cuando la extensión residual de $O$ sobre
$\Lambda$ es trivial (por ejemplo, si el cuerpo residual de $\Lambda$ es
algebraicamente cerrado), esto equivale a decir que $O$ mismo es isomorfo a un
álgebra de series formales de esta clase.
\footnote{La fe de erratas de 1962 precisa que esta definición de «liso sobre
$\Lambda$» solo es correcta cuando la extensión residual $k'/k$ es separable;
en el caso general, remite a SGA III, 1.1.}
Finalmente, si ya no se supone necesariamente que $O$ sea noetheriano, se dirá
también que $O$ es liso sobre $\Lambda$ si $O$ es isomorfo a un límite
proyectivo topológico de $\Lambda$-álgebras cociente que son noetherianas y
lisas sobre $\Lambda$ en el sentido anterior. Se generaliza inmediatamente al caso
en que $\Lambda$ y $O$ ya no se suponen locales. Dicho esto, se tiene la
proposición siguiente:

\result{PROPOSICIÓN}{5.2} Para que $O$ sea liso sobre $\Lambda$, es necesario
y suficiente que el funtor asociado $F$ transforme epimorfismos en
epimorfismos.

Esto significa que, para todo homomorfismo sobreyectivo
$A\longrightarrow A'$ en $\mathcal C$, $F(A)\longrightarrow F(A')$ es
sobreyectivo. Naturalmente, basta comprobar esta condición cuando $A$ es local
y, procediendo paso a paso, cuando el ideal de $A$ que es núcleo de
$A\longrightarrow A'$ está anulado por el ideal maximal de $A$. En la práctica,
esto reduce la cuestión a comprobar que cierta obstrucción, ligada a invariantes
infinitesimales

\src{377}
de la situación que da origen al funtor $F$, es nula, problema de naturaleza
cohomológica.

Para terminar, digamos unas palabras, en el contexto anterior, sobre los
\emph{anillos de definición}. Sea todavía $F$ un funtor de $\mathcal C$ en
$(\operatorname{Ens})$, pro-representable mediante una $\Lambda$-álgebra
topológica $O$. Entonces, para todo $A\in\mathcal C$ y todo $\xi\in F(A)$,
existe un menor subanillo $A'$ de $A$ tal que $\xi$ sea imagen de un elemento
$\xi'$ de $F(A')$ (que queda entonces determinado de manera única): en efecto,
basta interpretar $\xi$ mediante un homomorfismo de $O$ en $A$ y tomar como
$A'$ la imagen de $O$ por este último. Se dice entonces que $A'$ es el anillo
de definición del objeto $\xi\in F(A)$. Si $u:A\longrightarrow B$ es un
homomorfismo de álgebras y $\eta=F(u)(\xi)$, el anillo de definición de $\eta$
es la imagen por $u$ del anillo de definición de $\xi$. Cuando se parte de un
funtor $F$ de $\mathcal C$ en $(\operatorname{Ens})$, la existencia de los
anillos de definición y las propiedades que acabamos de dar son, con pequeñas
salvedades, equivalentes a que $F$ sea pro-representable; es decir, por lo
general distan mucho de ser triviales.

\fsection{B}{Los dos teoremas de existencia}

Conservemos las notaciones de A, párrafo 5, y partamos de un funtor covariante
\[
  F:\mathcal C\longrightarrow(\operatorname{Ens});
\]
se buscan criterios manejables para que $F$ sea pro-representable, es decir,
expresable mediante una $\Lambda$-álgebra $O$ como arriba. En virtud del
corolario de A, proposición 3.1, para ello es necesario y suficiente que $F$ sea
exacto por la izquierda. En el estado actual de la técnica de descenso (cf. las
cuestiones planteadas en [3], p. 9), este criterio no puede comprobarse
directamente bajo esta forma en los casos más importantes, y se necesitan
criterios aparentemente menos exigentes.

\result{TEOREMA}{1} Para que el funtor $F$ sea pro-representable, es necesario
y suficiente que satisfaga las dos condiciones siguientes:
\begin{enumerate}
\item[(i)] $F$ conmuta con los productos finitos;
\item[(ii)] para toda álgebra $A\in\mathcal C$ y todo homomorfismo
$A\longrightarrow A'$ en $\mathcal C$ tal que el diagrama
\[
  A\longrightarrow A'\rightrightarrows A'\otimes_A A'
\]
sea exacto ([3], A, definición 1.2), el diagrama transformado
\[
  F(A)\longrightarrow F(A')\rightrightarrows F(A'\otimes_A A')
\]

\src{378}
es exacto.
\end{enumerate}
Además, basta comprobar (ii) cuando $A$ es local y se está, además, en uno de
los dos casos siguientes:
\begin{enumerate}
\item[(a)] $A'$ es un módulo libre sobre $A$;
\item[(b)] el módulo cociente $A'/A$ es un $A$-módulo de longitud 1.
\end{enumerate}

La demostración de este teorema es bastante delicada y no puede esbozarse aquí.
Limitémonos a señalar que descansa esencialmente en un estudio de las relaciones
de equivalencia (en el sentido de las categorías) en el espectro de un álgebra
artiniana (estudio que plantea todavía varios problemas cuya solución parece
indispensable para el desarrollo ulterior de la teoría).

En las aplicaciones, la comprobación de la condición (i) es siempre trivial.
La de (ii) se descompone en dos: el caso en que $A'$ es un $A$-módulo libre
corresponde a la técnica de descenso por morfismos planos ([3], teoremas 1, 2 y
3) y no ofrece dificultad. Para tratar el caso (b), se utilizará el resultado
siguiente:

\result{TEOREMA}{2} Sean $A$ un anillo local artiniano de ideal maximal
$\mathfrak m$ y $A'$ una $A$-álgebra que contiene a $A$, tal que
$\mathfrak mA'\subset A$ y que
\[
  A\longrightarrow A'\rightrightarrows A'\otimes_A A'
\]
sea exacto (lo cual ocurre, en particular, si $A'/A$ es un $A$-módulo de
longitud 1). Sea $\mathcal F$ la categoría fibrada ([3], A, definición 1.1) de
los haces cuasicoherentes y planos sobre preesquemas variables. Entonces el
morfismo $\operatorname{Spec}(A')\longrightarrow\operatorname{Spec}(A)$ es un
morfismo de $\mathcal F$-descenso estricto ([3], A, definición 1.7).

En otros términos, el dato de un $A$-módulo plano $M$ es completamente
equivalente al dato de un $A'$-módulo plano $M'$, provisto de un
$A'\otimes_AA'$-isomorfismo de $M'\otimes_AA'$ sobre $A'\otimes_AM'$ que
satisface la condición de transitividad usual para un dato de descenso (loco
citato).

El teorema 2 se demuestra probando primero que
\[
  H^i(A'/A,G_a)=0\qquad\text{para }i\geq 1
\]
([3], A, 4(e)), pues la hipótesis $\mathfrak mA'\subset A$ permite reducirse
fácilmente al caso en que $A$ es un cuerpo $k$ (a saber, $A/\mathfrak m$). Se
aplican entonces las equivalencias señaladas en [3], página 16.

\fsection{C}{Aplicaciones a algunos casos particulares}

\subsection*{1. Observaciones generales sobre los funtores representados por preesquemas}
\addcontentsline{toc}{subsection}{1. Observaciones generales sobre los funtores representados por preesquemas}

Sea $S$ un preesquema localmente noetheriano. Se dice que un preesquema $X$
sobre $S$ es \src{379}localmente de tipo finito sobre $S$ si, para todo
$x\in X$ que se proyecta en $y\in S$,\footnote{El texto impreso dice aquí
« $y\in Y$ » y vuelve a emplear $Y$ al final del párrafo. Como el morfismo
considerado es $X\longrightarrow S$, se lee $S$ en ambas apariciones.} existen
un entorno afín de $y$, de anillo $A$, y un entorno afín de $x$ situado sobre
el anterior, de anillo $B$, tales que $B$ sea un álgebra de tipo finito sobre
$A$. Se encuentran muchos ejemplos importantes de preesquemas localmente de
tipo finito sobre $S$ que no son de tipo finito sobre $S$, como soluciones de
problemas universales clásicos; así, importa poder considerar el esquema de
Picard de una curva como la reunión de una infinidad de componentes conexas
(que no deben confundirse con la componente conexa del elemento neutro, es
decir, la «variedad de Picard»). En ocasiones resulta entonces conveniente
trabajar en la categoría $\mathcal C$ de los preesquemas localmente de tipo
finito sobre $S$, para examinar en ella la cuestión de la representabilidad de
un funtor contravariante $F$. El objetivo principal de estas exposiciones es
desarrollar una técnica general que permita reconocer si tal funtor $F$ es
representable y estudiar las propiedades del $S$-preesquema $X$ correspondiente
a partir de las de $F$. Señalemos de paso que, en este estudio, aparecen
ejemplos no patológicos de preesquemas sobre $S$ no separados sobre $S$, en
particular como «preesquemas de Picard» de excelentes $S$-esquemas; no deben,
por tanto, desterrarse de la geometría algebraica los preesquemas que no son
esquemas.

Sea $X$ un preesquema localmente de tipo finito sobre $S$, y sea
\[
  F(Y)=\operatorname{Hom}_S(Y,X)
\]
el funtor contravariante asociado. Se puede considerar la restricción $F_0$ de
$F$ a la subcategoría $\mathcal C_0$ de $\mathcal C$ formada por los
preesquemas $Y$ sobre $S$ que son artinianos y finitos sobre $S$: cuando
$S=\operatorname{Spec}(\Lambda)$, $\mathcal C_0$ es, pues, la categoría dual
de la categoría de las $\Lambda$-álgebras artinianas considerada en B. Si
$Y=\operatorname{Spec}(A)$, donde $A$ es un anillo local artiniano, $Y$ se
reduce a un solo punto $y$ situado sobre un punto cerrado $s$ de $S$, y un
$S$-homomorfismo de $Y$ en $X$ (es decir, un elemento de $F(Y)$) queda definido
por el dato de un punto $x\in X$ sobre $s$ y de un
$\mathcal O_{s,S}$-homomorfismo de $\mathcal O_{x,X}$ en $A$. Si existe tal
homomorfismo, entonces $x$ es necesariamente un punto cerrado de $X$ (pues su
cuerpo residual es algebraico sobre el de $s$). Esto muestra que la restricción
$F_0$ de $F$ a las «$\Lambda$-álgebras artinianas» es pro-representable, y está
representada por el álgebra topológica sobre $\Lambda$ cuyas componentes
locales son los completados $\widehat{\mathcal O}_{x,X}$ de los anillos locales
de $X$ en los puntos $x$ de $X$ que son cerrados y están situados sobre puntos
cerrados de $S$. Esto muestra que el solo conocimiento de $F_0$ proporciona
información precisa sobre la estructura de $X$ (a saber, la estructura de los
completados de sus anillos locales en los puntos indicados). Obsérvese que,
incluso cuando $S$

\src{380}
es el espectro de un cuerpo algebraicamente cerrado, solo mediante la
consideración sistemática de «variedades» $Y$ tales que $\mathcal O_Y$ pueda
tener elementos nilpotentes, en particular trabajando con los espectros de
anillos artinianos locales, se puede llegar a la «formulación correcta» de los
problemas universales clásicos y captar su aspecto «infinitesimal».

Cuando se parte de un funtor $F$ dado de antemano y se quiere decidir si es
representable, el estudio del funtor $F_0$ (con ayuda de los teoremas 1 y 2)
proporcionará indicaciones casi completas: o bien, como ocurre con frecuencia
(comprobando simplemente, por ejemplo, la naturaleza de los conjuntos
$F(I_K,\xi)$ y su comportamiento funtorial, cf. A), se constata que $F_0$ ya no
es pro-representable (lo que explica el fracaso de los intentos realizados
hasta ahora para definir de manera más o menos natural variedades de módulos
para la clasificación de los fibrados vectoriales de rango $>1$); o bien se
logra comprobar que $F_0$ es representable, pero los espacios vectoriales
$F(I_K,\xi)$ no son de dimensión finita, en cuyo caso hay que conformarse con
la solución «formal»; o bien, por último, $F_0$ está representado por un
producto de anillos locales completos noetherianos, lo que aporta indicios muy
sólidos de que el propio $F$ es representable y, unido a propiedades análogas
pero de naturaleza más global que pensamos desarrollar más adelante, bastará
sin duda para implicar que así sea efectivamente. Finalmente, se encuentran
problemas geométricos interesantes (véanse los párrafos 4 y 5 más abajo) en los
que solo se dispone del funtor $F_0$ (que no procede de ningún funtor «global»
$F$), y en los que habrá motivos para darse por satisfecho si se logra asociarle
un «esquema formal de módulos».

Para terminar estas generalidades, indiquemos cómo la teoría de esquemas
explica anomalías aparentes, como la superficie de Igusa $V$, cuya «variedad
de Picard» $P$ se reduce a un punto y para la cual, sin embargo, se tiene
\[
  H^1(V,\mathcal O_V)\ne 0 ;
\]
en este caso, $P$ es un grupo «puramente infinitesimal» que no se reduce al
grupo unidad, es decir, está definido por un álgebra local $O$ de rango finito
sobre el cuerpo base $k$ y dotado de una aplicación diagonal correspondiente a
la estructura multiplicativa de $P$; si $\mathfrak m$ es el ideal maximal de
$O$, el dual de $\mathfrak m/\mathfrak m^2$ es canónicamente isomorfo a
$H^1(V,\mathcal O_V)$ (cf. el párrafo 3 más abajo). Solo cuando el grupo de
Picard es un grupo algebraico en el sentido clásico, es decir, liso sobre el
cuerpo base $k$, la dimensión de $H^1(V,\mathcal O_V)$ (siempre igual a la de
$\mathfrak m/\mathfrak m^2$) es igual a la del grupo de Picard.

\subsection*{2. Esquemas $\underline{\operatorname{Hom}}_S(X,Y)$,
$\Pi_{X/S}Z$, $\underline{\operatorname{Aut}}_S(X)$, etc.}
\addcontentsline{toc}{subsection}{2. Esquemas Hom, Pi y Aut}

Sean $X,Y$ dos preesquemas sobre $S$; para todo preesquema $T$ sobre $S$,

\src{381}
sean $X_T=X\times_ST$, $Y_T=Y\times_ST$, y consideremos el conjunto
\[
  F(T)=\operatorname{Hom}_T(X_T,Y_T)
      =\operatorname{Hom}_S(X_T,Y)
      =\operatorname{Hom}_S(X\times_ST,Y)
\]
como funtor contravariante en $T$. Si es representable, se denota por
$\underline{\operatorname{Hom}}_S(X,Y)$ el preesquema sobre $S$ que lo
representa; se tiene, por tanto, un isomorfismo funtorial
\[
  \operatorname{Hom}_S
  \bigl(T,\underline{\operatorname{Hom}}_S(X,Y)\bigr)
  \xrightarrow{\sim}
  \operatorname{Hom}_S(T\times_SX,Y).
\]
Hay muchas variantes de este problema universal cuya solución se reduce a él:
el preesquema de los $S$-automorfismos de un $S$-preesquema $X$ (que será un
preesquema en grupos), el preesquema de los $S$-homomorfismos de un
$S$-preesquema en grupos en otro (que será un preesquema en grupos conmutativos
si el segundo esquema en grupos es conmutativo), etc. También se puede
generalizar la definición de $\underline{\operatorname{Hom}}_S(X,Y)$
considerando un preesquema $Z$ sobre el preesquema $X$ sobre $S$, y el funtor
\[
  F(T)=\operatorname{Hom}_{X_T}(X_T,Z_T)
\]
(conjunto de las «secciones» de $Z_T$ sobre $X_T$); si este funtor es
representable, el $S$-preesquema que lo representa se denotará por
$\Pi_{X/S}Z$; se tendrá, pues, por definición, un isomorfismo funtorial
\[
  \operatorname{Hom}_S(T,\Pi_{X/S}Z)
  =\operatorname{Hom}_{X_T}(X_T,Z_T).
\]
Tomando $Z=Y\times_SX$, se recupera
$\underline{\operatorname{Hom}}_S(X,Y)$. De estas definiciones resulta para los
nuevos preesquemas así introducidos un formulario tan trivial como útil, que no
daremos aquí (puesto que es válido en toda categoría en la que existen
productos y productos fibrados). Más seria es la cuestión de la existencia de
esquemas del tipo $\underline{\operatorname{Hom}}_S(X,Y)$. Se constata primero
que, para $X$ fijo, $\underline{\operatorname{Hom}}_S(X,Y)$ solo puede existir
para todo $Y$ sobre $S$ (o $\Pi_{X/S}Z$ para todo $Z$ sobre $X$) si $X$ es plano
sobre $S$. Además, se advierte que solo es razonable esperar la existencia de
una solución, para $Y$ suficientemente general, si $X$ es además propio sobre
$S$. Por otra parte, parece que estas condiciones son suficientes para la
existencia de $\underline{\operatorname{Hom}}_S(X,Y)$ y $\Pi_{X/S}Z$, a
condición de imponer, cuando proceda, una hipótesis del tipo «cuasiproyectivo»
sobre $Y/S$, respectivamente sobre $Z/X$; esto puede comprobarse, en cualquier
caso, en bastantes situaciones (por ejemplo, cuando $Y$ es afín sobre $S$; o,
mediante construcciones directas elementales, cuando $X$ es finito sobre $S$).
He aquí lo que dan los teoremas 1 y 2:

\src{382}
\result{PROPOSICIÓN}{2.1} Sean $\Lambda$ un anillo noetheriano, $X$ e $Y$ dos
preesquemas cualesquiera sobre $\Lambda$, y consideremos el funtor
\[
  F(A)=\operatorname{Hom}_A(X_A,Y_A)
\]
en la categoría $\mathcal C_0$ de las $\Lambda$-álgebras artinianas. Si $X$ es
plano sobre $\Lambda$, este funtor es pro-representable.

Además, se comprueba que para todo $A\in\mathcal C_0$ y $\xi\in F(A)$ se tiene
un isomorfismo canónico
\[
  F(I_A,\xi)=H^0\!\left(
    X_A,
    \underline{\operatorname{Hom}}_{\mathcal O_{X_A}}
    \bigl(\xi^*(\Omega^1_{Y_A/A}),\mathcal O_{X_A}\bigr)
  \right),
\]
donde $\Omega^1_{Y_A/A}$ es el haz de $1$-diferenciales de Kähler de $Y_A$
respecto de $A$. Tomando como $A$ un cuerpo, se obtiene, utilizando A.5.1 y el
teorema de finitud de [2], el corolario siguiente:

\result{COROLARIO}{} Supongamos que $X$ es plano y propio sobre $S$, y que $Y$
es localmente de tipo finito sobre $S$. Entonces $F$ es pro-representable y las
componentes locales del álgebra topológica correspondiente sobre $\Lambda$ son
anillos noetherianos.

\result{OBSERVACIONES}{} Los problemas considerados en este apartado, y muchos
otros, se estudiaban habitualmente, en el marco de la geometría algebraica
clásica, mediante las «coordenadas de Chow» de los ciclos en el espacio
proyectivo, que permitían considerar dichos ciclos como puntos de variedades
proyectivas adecuadas. Este procedimiento, y en general el uso de las
coordenadas de Chow, parece irremediablemente insuficiente desde el punto de
vista de los esquemas, pues destruye los elementos nilpotentes en las
variedades de parámetros y, en particular, no se presta a un estudio
satisfactorio de las variaciones infinitesimales de ciclos (por no hablar de su
naturaleza no intrínseca, ligada al espacio proyectivo). El lenguaje de las
coordenadas de Chow era, por desgracia, el único utilizado por muchos
geómetras algebraicos para estudiar familias de variedades o familias de
ciclos, lo que parece haber constituido un obstáculo serio para aclarar estas
nociones, pese a su indudable interés técnico (probablemente provisional). Si
se quiere obtener el análogo de las variedades de Chow en la teoría de
esquemas, se llega al problema universal siguiente: sea $X$ un preesquema sobre
$S$; para todo preesquema $T$ sobre $S$, se considera el conjunto $F(T)$ de los
subpreesquemas cerrados de $X_T=X\times_ST$ que son planos sobre $T$. Se intenta
representar este funtor en $T$ mediante un preesquema sobre $S$. Más
generalmente, se puede dar un haz cuasicoherente $G$ sobre $X$ y tomar como
$F(T)$ el conjunto de los haces

\src{383}
cuociente de $G_T$ que son planos sobre $T$. Parece que el problema tiene una
solución, mediante un esquema $C$ localmente de tipo finito sobre $S$, cuando
$X$ es propio sobre el preesquema localmente noetheriano $S$ y $G$ es además
coherente. En cualquier caso, suponiendo únicamente que $S$ es localmente
noetheriano, la restricción de $F$ a las «$S$-álgebras artinianas» es
pro-representable; y si, además, $X$ es propio sobre $S$ y $G$ es coherente, las
componentes locales del anillo topológico correspondiente $O$ son
noetherianas. Por supuesto, aun después de demostrar la existencia del
«esquema de Chow» de $X$ sobre $S$, quedará por
encontrar una descomposición en abiertos disjuntos $C_i$ (correspondientes a
la fijación de invariantes continuos, como el grado y la dimensión de los
ciclos que se hacen variar) que sean de tipo finito (y no solo localmente de
tipo finito) sobre $S$, precisar las relaciones de este esquema con las
variedades de Chow clásicas y precisar cuándo un $C_i$ es incluso proyectivo,
o al menos cuasiproyectivo, sobre $S$.\footnote{La fe de erratas de 1962
precisa que el caso proyectivo queda completamente resuelto por los esquemas de
Hilbert (TDTE IV), mientras que ejemplos de Nagata y Hironaka muestran que, sin
una hipótesis de proyectividad, estos funtores pueden no ser representables,
incluso para las subvariedades de dimensión cero de una variedad completa no
singular de dimensión tres; este fenómeno está ligado a la posible
inexistencia del cuadrado simétrico.}

\subsection*{3. Esquemas de Picard}
\addcontentsline{toc}{subsection}{3. Esquemas de Picard}

Sea $f:X\longrightarrow S$ un $S$-preesquema,\footnote{Para un estudio más
completo, la fe de erratas de 1962 remite a TDTE V.} consideremos el haz
multiplicativo $\mathcal O_X^*$ de las unidades del haz estructural de $X$, y
el grupo
\[
  P(X/S)=H^0\bigl(S,R^1f_*(\mathcal O_X^*)\bigr),
\]
llamado grupo de Picard relativo de $X/S$. Un elemento de este grupo queda,
pues, definido dando una cubierta abierta $(U_i)$ de $S$ y un haz invertible
$L_i$ sobre cada $f^{-1}(U_i)$, de modo que, para todo par $(i,j)$,
$L_i|f^{-1}(U_i\cap U_j)$ sea isomorfo a $L_j|f^{-1}(U_i\cap U_j)$, al menos
localmente sobre $U_i\cap U_j$ (es decir, estos dos haces sean «equivalentes»
en el sentido de [3], B, 4). Si $X/S$ admite una sección, entonces $P(X/S)$ no
es sino el conjunto de las clases de haces invertibles sobre $X/S$ módulo
«equivalencia» (loco citato). Pongamos ahora, para todo $T$ sobre $S$,
\[
  F(T)=P(X_T/T) ;
\]
se obtiene un funtor contravariante en $T$, que podrá llamarse funtor de Picard de
$X/S$; si este funtor es representable, el preesquema sobre $S$ que lo
representa se llamará preesquema de Picard de $X/S$ y se denotará por
$\underline P(X/S)$. Se tendrá, pues, un isomorfismo de funtores
\[
  \operatorname{Hom}_S\bigl(T,\underline P(X/S)\bigr)
  \xrightarrow{\sim}P(X_T/T).
\]
La formación de preesquemas de Picard es compatible con la extensión de la
base; en particular, los preesquemas de Picard de las fibras de $X$ sobre $S$
(que son preesquemas sobre los cuerpos residuales $K(s)$ de los $s\in S$) son
las fibras de $\underline P(X/S)$. Por supuesto, como $P(X_T/T)=F(T)$ es un
grupo conmutativo, los preesquemas de Picard son preesquemas en grupos.
Obsérvese también que las jacobianas generalizadas de Rosenlicht no son sino
las componentes conexas de la unidad en los esquemas de Picard de curvas
completas que pueden tener singularidades, lo que debería hacer evidentes la
mayoría de sus propiedades (una vez demostrada la existencia).

\result{OBSERVACIÓN}{} La definición adoptada aquí solo es razonable cuando
todo punto de $S$\footnote{El texto impreso dice «$Y$»; aquí no se ha definido
ninguna base $Y$, y la condición se refiere a la base $S$.} admite un entorno
abierto $U$ sobre el cual $X$ admite una sección. En el caso general, hay que
modificar ligeramente la definición del funtor de Picard para poder obtener
también un teorema de existencia.

Aquí, las condiciones plausibles para la existencia de un preesquema de Picard
son las siguientes: $X$ es propio y plano sobre $S$,
$f_*(\mathcal O_X)=\mathcal O_S$, y $X$ admite localmente una sección sobre
$S$. Esta condición aparece de manera natural al aplicar la técnica de
descenso, para eliminar los automorfismos de un haz invertible $L$ sobre $X$
dotándolo de una sección marcada sobre la sección $s$ ([3], B, 4). Se obtiene,
en particular:

\src{384}
\result{PROPOSICIÓN}{3.1} Supongamos que $X$ es plano sobre
$S=\operatorname{Spec}(\Lambda)$, con $\Lambda$ noetheriano, y que, para todo
$T$ de tipo finito sobre $S$, se tiene
\[
  f_{T*}(\mathcal O_{X_T})=\mathcal O_T
\]
(si $f$ es propio y separable,\footnote{La fe de erratas de 1962 prescribe
explícitamente la hipótesis «$f$ propio y separable».} o si $S$ es el espectro
de un cuerpo, resulta de Künneth que la última condición equivale a
$f_*(\mathcal O_X)=\mathcal O_S$). Entonces el funtor de Picard de $X/S$ en la
categoría de las $\Lambda$-álgebras artinianas es pro-representable.

Además, aquí se tendrá
\[
  F(I_A,\xi)=H^1(X_A,\mathcal O_{X_A}),
\]
en particular:

\result{COROLARIO}{} Si $X$ es propio sobre $S$, entonces las componentes
locales de la $\Lambda$-álgebra topológica $O$ correspondiente al funtor de
Picard son noetherianas.

\result{OBSERVACIONES}{} Se pueden generalizar las definiciones y los resultados
de este apartado a la clasificación de los fibrados principales sobre $X$, con
grupo estructural un esquema en grupos $G$ sobre $S$ que sea afín, plano sobre
$S$ y conmutativo. Cuando $G$ no es conmutativo, y por tanto el fibrado en
grupos adjunto de un fibrado principal (cuyas secciones son los automorfismos
del fibrado principal) deja de ser trivial, la proposición 3.1 ya no es válida
tal cual. No obstante, se puede modificar el problema universal de manera que
siga teniendo solución (al menos, por ahora, en geometría formal). La regla de
oro que debe recordarse, en el contexto de este \src{385}apartado y de los
siguientes, y siempre que se busquen «esquemas de módulos» para clases de
objetos definidos solo salvo isomorfismo, es siempre la siguiente: eliminar
los posibles automorfismos de los objetos que se quieren clasificar mediante
la introducción, si es necesario, de estructuras auxiliares (puntos o
elementos de secciones marcadas, fijación de formas diferenciales, etc.), que
se elegirán lo bastante inocuas como para no modificar sustancialmente el
problema inicial.

\subsection*{4. Módulos formales de una variedad}
\addcontentsline{toc}{subsection}{4. Módulos formales de una variedad}

Sea $\Lambda$ un anillo local noetheriano de cuerpo residual $k$ (con mayor
frecuencia, $\Lambda$ será igual a $k$, o a un $p$-anillo de Cohen), y sea
$X_0$ un preesquema sobre $k$. Para toda $\Lambda$-álgebra artiniana local $A$,
se considera el conjunto $F(A)$ de las clases, salvo isomorfismo, de
$A$-preesquemas $X$ planos sobre $A$, dotados de un isomorfismo
\[
  (*)\qquad
  X\otimes_A\kappa(A)\xleftarrow{\sim}
  X_0\otimes_k\kappa(A),
\]
donde $\kappa(A)$ es el cuerpo residual de $A$; por supuesto, los isomorfismos
entre tales $A$-preesquemas planos deben respetar el isomorfismo precedente
incluido en la estructura. Si $A$ es una $\Lambda$-álgebra artiniana no
necesariamente local, de componentes locales $A_i$, se toma como $F(A)$ el
producto de los $F(A_i)$. De este modo, $F$ se convierte en un funtor
multiplicativo en $A$, que podrá llamarse funtor de módulos para $X_0$ (y
$\Lambda$). Si este funtor es representable, le corresponde una
$\Lambda$-álgebra topológica $O$ local, de cuerpo residual $k$, y el espectro
formal de $O$ se llamará esquema formal de módulos para $X_0$ (y $\Lambda$);
cf. [2] para algunos detalles sobre esta noción.

Aquí, para aplicar la técnica de descenso, los automorfismos «finitos» de $X_0$
son inofensivos, pues no influyen en la existencia de automorfismos (en el
sentido precisado más arriba) de los $A$-preesquemas $X$; cuando $A$ no se
reduce a un cuerpo, la condición necesaria y suficiente para que $X$ no tenga
ningún $A$-automorfismo no trivial es que se tenga
\[
  H^0(X_0,\Theta_{X_0/k})=0,
\]
donde $\Theta_{X_0/k}$ denota el haz de las $k$-derivaciones (= haz tangente)
de $X_0$. Por otra parte, se obtiene fácilmente (al menos si $X_0$ es liso
sobre $k$)
\[
  F(I_A,\xi)=H^1(X_A,\Theta_{X_A/A}).
\]
Se concluye entonces, como de costumbre:

\src{386}
\result{PROPOSICIÓN}{4.1} Supongamos que
$H^0(X_0,\Theta_{X_0/k})=0$. Entonces existe el esquema formal de módulos para
$X_0$. Si, además, $X_0$ es propio sobre $k$, el esquema formal de módulos es
noetheriano.

\result{OBSERVACIONES}{} 1\textsuperscript{o} Cuando no se supone que $X_0$ sea
liso sobre $k$, $F(I_A,\xi)$ se identifica con un sub-$A$-módulo de
\[
  \operatorname{Ext}^1_{\mathcal O_{P_A}}
  \bigl(P_A;\mathcal I_{X_A},\mathcal O_{X_A}\bigr),
\]
donde se pone $P_A=X_A\times_\Lambda X_A$, donde $\mathcal O_{X_A}$ se
considera como un haz coherente sobre $P_A$ mediante el morfismo diagonal
$X_A\longrightarrow P_A$, y donde $\mathcal I_{X_A}$ denota el haz coherente
de ideales sobre $P_A$ definido por el morfismo diagonal. Más precisamente,
una globalización sencilla de la teoría de Hochschild muestra que el
$\operatorname{Ext}^1$ anterior se identifica con el conjunto de las clases,
salvo isomorfismo, de haces de $I_A$-álgebras planas $\mathcal O$ sobre $X_A$,
dotados de un isomorfismo de aumentación
\[
  \mathcal O\otimes_{I_A}A\longrightarrow\mathcal O_{X_A}
\]
(se recuerda que se pone $I_A=A[t]/(t^2)$). El submódulo $F(I_A,\xi)$ es el
que corresponde a los haces de álgebras conmutativas. Las hipótesis de lisura
no son, pues, esenciales en la teoría de módulos, contrariamente a lo que [2]
daba a entender.

2\textsuperscript{o} Recordemos (loco citato) que, en particular, toda curva
algebraica $X_0$ lisa y propia sobre $k$ admite un esquema formal de módulos,
liso sobre $\Lambda$, de dimensión relativa $3g-3$ si el género $g$ es
$\geq2$, y $g$ si $g=0$ o $1$. Estos dos últimos casos ya no entran directamente
en la proposición 4.1. Sin embargo, en el caso de las curvas elípticas
($g=1$), podrían reducirse a ella mediante las consideraciones que siguen.

Por supuesto, se puede variar ad libitum la proposición 4.1 considerando
sistemas de esquemas sobre $k$ dotados de diversas estructuras. Supongamos,
por ejemplo, que $X_0$ sea un esquema abeliano sobre $k$, con punto origen
marcado (es decir, que $X_0$ se considere como esquema en grupos sobre $k$), y
sea $F(A)$ el conjunto de las clases, salvo isomorfismo, de esquemas abelianos
sobre $A$ (es decir, de esquemas en grupos propios y lisos sobre $A$) dotados
de un isomorfismo de esquemas abelianos $(*)$. Se comprueba que la imposición
de una estructura multiplicativa (e incluso de una mera «sección unidad»)
elimina los automorfismos infinitesimales y que, por consiguiente, existe un
esquema formal de módulos, correspondiente a un anillo local completo
noetheriano $O$. Por otra parte, puede demostrarse que, si $X$ es un esquema
propio y liso con fibras «absolutamente conexas» sobre un preesquema localmente
noetheriano $S$, toda estructura multiplicativa sobre $X$ que admita una
sección unidad es necesariamente asociativa y conmutativa (siempre que lo sea,
al menos, sobre una fibra y que $S$ sea conexo), y además queda determinada de
manera única por el conocimiento

\src{387}
de la sección unidad. Además, supongamos que $S$ es el espectro de un anillo
local artiniano $A$ de cuerpo residual $k$, que $X$ es propio sobre $A$ y está
dotado de una sección $s$, y que, por último, $X\otimes_Ak$ está dotado de una
estructura de esquema abeliano sobre $k$ cuyo elemento neutro es el punto de
$X\otimes_Ak$ correspondiente a $s$; un cálculo sencillo de obstrucciones,
unido a un razonamiento debido a Serre, permite demostrar que existe
efectivamente sobre $X$ una estructura multiplicativa que admite la sección
$s$ como sección unidad. (A partir de ahí, utilizando el «teorema de
existencia» de [2] para pasar al caso en que $A$ es noetheriano local completo,
y luego la técnica de descenso de [3] para el caso general, se puede demostrar
el enunciado análogo para todo $S$ localmente noetheriano conexo.) Esto prueba
que el funtor $F(A)$ considerado aquí es isomorfo al funtor análogo definido al
principio de este apartado prescindiendo de la estructura multiplicativa sobre
$X_0$. Se sigue, en particular, que, si $\mathfrak m$ es el ideal maximal de
$O$, $\mathfrak m/\mathfrak m^2$ es canónicamente isomorfo al dual de
$H^1(X_0,\Theta_{X_0/k})$ y tiene, por tanto, dimensión $n^2$, donde
$n=\dim X_0$. Sería muy interesante determinar si $O$ es liso sobre $\Lambda$,
es decir, isomorfo a un álgebra de series formales en $n^2$ indeterminadas
sobre $\Lambda$. El párrafo 1, proposición 5.2, permite dar una formulación
equivalente de este problema como problema de existencia de esquemas abelianos
que se reduzcan según un esquema abeliano dado. En cualquier caso, por métodos
trascendentes se ve que la respuesta es afirmativa si $k$ es de característica
cero. En característica $p\ne0$, basta evidentemente limitarse al caso en que
$\Lambda$ es el anillo de vectores de Witt construido sobre un cuerpo
algebraicamente cerrado $k$. Tal vez haya llegado el momento de que el «funtor
de Greenberg» demuestre su valía\ldots\footnote{La fe de erratas de 1962
anuncia que el esquema formal de módulos de una variedad abeliana sobre un
cuerpo es liso sobre el anillo de Witt: todo esquema abeliano sobre un cociente
artiniano local se levanta al anillo del cual es cociente. La demostración
utiliza la varianza de la clase de obstrucción al levantamiento introducida en
la exposición 182, p. 12. También recuerda las construcciones de Mumford para
las curvas de género $g$ y los esquemas abelianos polarizados (cf. SHMT).}

\subsection*{5. Extensión de los revestimientos}
\addcontentsline{toc}{subsection}{5. Extensión de los revestimientos}

Sean $\mathfrak X$ un preesquema formal noetheriano [2] y $U$ una parte
abierta de $\mathfrak X$, definida localmente por la «no anulación» de una
sección de $\mathcal O_{\mathfrak X}$ que no sea divisor de cero no solo sobre
$\mathfrak X$, sino sobre todo $X_n$,\footnote{Precisión impuesta por la fe de
erratas de 1962.} y que sea, por tanto, lo bastante grande para que toda sección
de $\mathcal O_{\mathfrak X}$ sobre un abierto $V$ que sea nula sobre $U\cap V$
sea nula. Sea $\mathcal J$ un «ideal de definición» para $\mathfrak X$, y sea
\[
  X_n=(\mathfrak X,\mathcal O_{\mathfrak X}/\mathcal J^{n+1}),
\]
que es, por tanto, un preesquema noetheriano ordinario. Si entonces
$\mathfrak X',\mathfrak X''$ son dos revestimientos planos de $\mathfrak X$ (es
decir, dos preesquemas sobre $\mathfrak X$ definidos por haces de álgebras que
son coherentes y localmente libres como haces de módulos), no ramificados sobre
$U$, la aplicación evidente
\[
  \operatorname{Hom}_{\mathfrak X}(\mathfrak X',\mathfrak X'')
  \longrightarrow
  \operatorname{Hom}_{X_0}(X'_0,X''_0)
\]
es inyectiva; en particular, un automorfismo de $\mathfrak X'$ que induce la
identidad sobre $X'_0$ es la identidad. Esto permite aplicar a la situación la
técnica de descenso. Partamos, en particular, de un revestimiento plano $X'_0$
de $X_0$, no ramificado

\src{388}
sobre $U_0$; sea $G(\mathfrak X)$ el conjunto de las clases, salvo isomorfismo
(que induzca la identidad sobre $X'_0$), de revestimientos planos $\mathfrak X'$
de $\mathfrak X$ que inducen $X'_0$ sobre $X_0$ (y son, por tanto,
necesariamente no ramificados sobre $U$). Se define de igual modo $G(V)$ para
toda parte abierta $V$ de $\mathfrak X$, y más generalmente $G(\mathfrak Y)$
para todo preesquema formal $\mathfrak Y$ sobre $\mathfrak X$. Dicho esto, los
resultados de [2] y [3] implican en primer lugar los resultados siguientes:

\noindent a. Cuando $V$ es un abierto variable de $\mathfrak X$, los $G(V)$
forman un haz sobre $\mathfrak X$, sea $G_{\mathfrak X}=G$. La restricción de
este haz a $U$ es el haz constante cuyas fibras se reducen a un elemento.

Más generalmente, la determinación de las fibras de $G_{\mathfrak X}$ es una
cuestión de anillos locales completos; más precisamente:

\noindent b. Para todo $x\in\mathfrak X$, se tiene
\[
  G_x=G\bigl(\operatorname{Spec}(\mathcal O_{\mathfrak X,x})\bigr)
  \subset
  G\bigl(\operatorname{Spec}(\widehat{\mathcal O}_{\mathfrak X,x})\bigr),
\]
(= clases salvo isomorfismo de álgebras finitas y libres $B$ sobre
$\widehat{\mathcal O}_{\mathfrak X,x}$ dotadas de un isomorfismo de
$B\otimes_{\widehat{\mathcal O}_{\mathfrak X,x}}\mathcal O_{X_0,x}$ con
$\widehat{\mathcal O}'_{0x}$, donde $\mathcal O'_0$ es el haz de álgebras sobre
$X_0$ que define $X'_0$).\footnote{La fórmula impresa identifica directamente
$G_x$ con el término completado y omite el sombrero sobre
$\mathcal O'_{0x}$. La inclusión y el completado son los prescritos por la fe
de erratas de 1962.}

\noindent c. Se tiene un isomorfismo canónico
\[
  G_{\mathfrak X}=\varprojlim_nG_{X_n} ;
\]
en otros términos, para todo abierto $V$ de $\mathfrak X$, se tiene
\[
  G(V)=\varprojlim_nG(V_n).
\]

\noindent d. Supongamos que $\mathfrak X$ se obtiene de un esquema ordinario
propio $X$ sobre un anillo local noetheriano completo $\Lambda$ que posee un
ideal de definición $\mathfrak m$, tomando el completado $\mathcal J$-ádico de
$\mathcal O_X$, donde $\mathcal J=\mathfrak m\mathcal O_X$. Entonces
$G(\mathfrak X)$ es canónicamente isomorfo al conjunto de las clases de
revestimientos planos del esquema ordinario $X$ que «se reducen a $X'_0$».

De manera gráfica, puede decirse que (a) y (b) establecen las relaciones
fundamentales entre el aspecto local y el aspecto global del problema, (c) da
las relaciones entre el aspecto «finito» y el aspecto «infinitesimal» y, por
último, (d) recuerda (bajo las condiciones precisadas) la identidad entre el
aspecto «formal» y el aspecto «algebraico».

Supongamos ahora que $\mathfrak X$ está definido sobre un anillo local
noetheriano completo $\Lambda$, con
$\mathcal J=\mathfrak m\mathcal O_{\mathfrak X}$ (por tanto, $X_0$ es un
preesquema sobre $\Lambda/\mathfrak m$). Para toda álgebra $A$ finita sobre
$\Lambda$, pongamos
\[
  F(A)=G(\mathfrak X\times_\Lambda A),
\]
donde el marcado se toma después del cambio de base al álgebra reducida $A_{\mathrm{red}}$, y no literalmente a $A/\mathfrak mA$. Es un funtor covariante en $A$, con valores en la categoría de los conjuntos,
y por (c) este funtor queda completamente determinado si se lo conoce para $A$
artiniana; y decir que este funtor es pro-representable, es decir, de la forma
\[
  F(A)=\operatorname{Hom}_{\Lambda\text{-algèbres top.}}(O,A),
\]

\src{389}
donde $O$ es un álgebra topológica sobre $\Lambda$ del tipo considerado en A,
párrafo 5, equivale a decir que lo mismo ocurre al restringirse a los $A$
artinianos. La conjunción de los teoremas 1 y 2 implica entonces efectivamente:

\result{PROPOSICIÓN}{5.1} Siempre que toda familia artiniana clasificada por el funtor sea plana sobre su álgebra de coeficientes y que el abierto escogido permanezca esquemáticamente denso después de los cambios de base pertinentes, el funtor precedente es pro-representable.

Por supuesto, según (a), si $U=\mathfrak X$, $G(\mathfrak Y)$ se reduce a un
solo elemento para todo $\mathfrak Y$ sobre $\mathfrak X$, y el funtor $F$
tiene entonces escaso interés (se tendrá $O=\Lambda$). Parece que, en casi
cualquier otro caso, el anillo local topológico $O$ no es noetheriano. No
obstante, su existencia pone de manifiesto, de manera sorprendente, la
naturaleza «continua» del conjunto $G(\mathfrak X)$ de las soluciones (que
corresponde intuitivamente al hecho de que hay una elección «continua» en el
modo en que la ramificación se propaga al extender $X'_0$). Compárese este
resultado con el punto de vista de J.-P. Serre [4] en la teoría del cuerpo de
clases local, que pone asimismo de relieve el carácter continuo del grupo de
Galois topológico de la extensión abeliana maximal de un cuerpo local
«geométrico»; allí, el grupo dual (en el sentido de Pontryagin) aparece como un
límite directo de grupos algebraicos (o al menos cuasialgebraicos); también
aquí, la clasificación de las extensiones viene dada por «variedades» de
dimensión infinita. Por otra parte, en lo anterior se puede tomar como
$\mathfrak X$ el espectro formal de un anillo local completo (del cual
$\Lambda$ será, por ejemplo, un subanillo de Cohen), y cabe esperar que los
resultados de este apartado puedan utilizarse en el estudio de las extensiones
de un anillo local completo de dimensión $>1$. Tanto en el caso local como en
el global, quizá permitan formular relaciones precisas entre los fenómenos de
ramificación superior y fenómenos en característica cero (abordables por
métodos trascendentes). En cualquier caso, el análisis preliminar de la
proposición 5.1 es el que permite extender al caso «moderadamente ramificado»
los métodos expuestos en [2] para estudiar el grupo fundamental y resolver por
métodos trascendentes el «problema de los tres puntos».

Para terminar, señalemos que la situación se simplifica cuando $X_0$ tiene dimensión 1: entonces, en virtud de (a), complementado por el pegado formal en los esquemas locales completados perforados,
\[
  G(\mathfrak X)\simeq
  \prod_iG\bigl(\operatorname{Spec}(\widehat{\mathcal O}_{\mathfrak X,x_i})\bigr),
\]
donde los $x_i$ son los puntos de $X_0-U$: pueden darse arbitrariamente las
extensiones «locales» en los puntos de ramificación. Además, la normalidad de $X_0$ por sí sola no implica que el esquema formal de módulos de 5.1 sea liso sobre $\operatorname{Spec}(\Lambda)$; la última afirmación impresa es falsa en general.

\src{390}
\begin{center}
  \textbf{BIBLIOGRAFÍA}
\end{center}

\begin{description}
  \item[{[1]}] Grothendieck (Alexander). --- \emph{Sur quelques points
    d'algèbre homologique}, Tohoku Math. J., t. 9, 1957, p. 119--221.
  \item[{[2]}] Grothendieck (Alexander). --- \emph{Géométrie formelle et
    géométrie algébrique}, Séminaire Bourbaki, t. 11, 1958/59, fasc. 3,
    n\textsuperscript{o} 182.
  \item[{[3]}] Grothendieck (Alexander). --- \emph{Technique de descente et
    théorèmes d'existence en géométrie algébrique, I : Généralités, Descente
    par morphismes fidèlement plats}, Séminaire Bourbaki, t. 12, 1959/60,
    fasc. 1, n\textsuperscript{o} 190.
  \item[{[4]}] Serre (Jean-Pierre). --- \emph{Corps locaux et isogénies},
    Séminaire Bourbaki, t. 11, 1958/59, fasc. 3, n\textsuperscript{o} 185.
\end{description}
